\secnumberlesssection{INTRODUCCIÓN}

La gestión logística constituye hoy un pilar de la economía global y un factor determinante de la competitividad empresarial. Dentro de este ámbito, la planificación eficiente de rutas de transporte es una necesidad transversal a industrias como la gestión de residuos, la distribución de alimentos o los servicios de mantenimiento. De este contexto surge el Problema de Ruteo de Vehículos Periódico (PVRP), que busca diseñar rutas óptimas para una flota que debe atender a un conjunto de clientes con una frecuencia determinada a lo largo de un horizonte de planificación de varios días. Tradicionalmente, esta familia de problemas se ha abordado mediante heurísticas y metaheurísticas; sin embargo, su dependencia de reglas fijas y de una calibración manual meticulosa limita su capacidad de adaptación frente a los entornos dinámicos que caracterizan a la logística moderna.

El PVRP es un problema de optimización combinatoria clasificado como NP-hard. Su dificultad no proviene únicamente del ruteo espacial, sino de la doble decisión que impone: determinar simultáneamente en qué días se visita a cada cliente y qué rutas se ejecutan en cada jornada. Esta combinación de decisiones temporales y espaciales expande el espacio de soluciones de forma exponencial y vuelve inviables a los métodos exactos en instancias de tamaño real. Frente a este panorama, el Aprendizaje por Refuerzo (RL) ha emergido como una alternativa prometedora, ya que permite que un agente aprenda políticas de decisión directamente desde la interacción con el entorno, en lugar de depender de reglas diseñadas manualmente. No obstante, su aplicación a variantes con dimensión temporal como el PVRP sigue siendo un área incipiente, sin modelos consolidados que aborden de forma conjunta la asignación de días y la construcción de rutas periódicas.

Esta memoria aborda dicha brecha desarrollando y evaluando un modelo basado en RL para la resolución del PVRP, bajo un enfoque data-driven. La propuesta reformula el problema como un proceso de decisión de Markov (MDP), en el que un agente construye la solución de manera secuencial, seleccionando los nodos de la ruta uno a uno y decidiendo en cada paso si continuar la ruta actual o cerrarla. Bajo este diseño, la asignación de los días de visita no se impone mediante reglas externas, sino que emerge de forma implícita del contexto de cada decisión. Sobre esta formulación se implementa un agente entrenado con el algoritmo Maskable PPO, cuyo desempeño se contrasta con dos métodos de referencia: una heurística voraz (Greedy) y la metaheurística Variable Neighborhood Search (VNS). La validación se organiza en torno a tres ejes, calidad, robustez y generalización, evaluados sobre instancias del conjunto NEO~\cite{NEO2013} agrupadas en tres familias estructurales y tres escalas de clientes, con un análisis estadístico multi-semilla que respalda la reproducibilidad de los resultados.

El presente documento se organiza en cinco capítulos. El Capítulo~\ref{sec:def} define el problema, su contexto, complejidad y formulación matemática, junto con los objetivos del trabajo. El Capítulo 2 presenta el marco conceptual y el estado del arte, cubriendo la optimización combinatoria, el aprendizaje por refuerzo y su aplicación a problemas de ruteo. El Capítulo 3 desarrolla la propuesta de solución, detallando la formulación del PVRP como MDP, la arquitectura del agente y los métodos de comparación y evaluación. El Capítulo 4 expone la validación experimental, con los resultados por escala, la identificación de los límites del método y el análisis de generalización. Finalmente, el capítulo 5 sintetiza los hallazgos, discute las limitaciones y contribuciones del trabajo, y plantea las líneas de investigación futuras.
